\documentclass[11pt,a4paper]{article}
\usepackage[margin=2.5cm]{geometry}
\usepackage{amsmath,amssymb,amsthm,mathtools}
\usepackage{booktabs,array,enumitem}
\usepackage[hidelinks]{hyperref}
\usepackage{microtype}

\newtheorem{theorem}{Theorem}
\newtheorem{proposition}[theorem]{Proposition}
\newtheorem{lemma}[theorem]{Lemma}
\newtheorem{corollary}[theorem]{Corollary}
\theoremstyle{definition}
\newtheorem{definition}[theorem]{Definition}
\newtheorem{example}[theorem]{Example}
\newtheorem{assumption}{Standing assumption}
\renewcommand{\theassumption}{S\arabic{assumption}}
\theoremstyle{remark}
\newtheorem{remark}[theorem]{Remark}

\newcommand{\E}{\mathbb{E}}
\newcommand{\Vend}{V_{\mathrm{end}}}
\newcommand{\pp}{\hat p}
\newcommand{\kt}{\tilde\kappa}
\newcommand{\Pt}{\tilde{\mathbb P}}
\newcommand{\Et}{\tilde{\mathbb E}}
\newcommand{\pos}[1]{\left(#1\right)^{+}}

\title{The conjectures in \emph{A Bayesian Algorithm for Optimal Dating Expenditure}:\\ proofs, counterexamples and corrected statements}
\author{Companion note}
\date{30 October 2026}

\begin{document}
\maketitle

\begin{abstract}
The proposal states five conjectures (Conjectures 1--5) and leaves one question open: whether the value of information can be negative once refusals sting (Section~4.4). None of the five conjectures holds as stated in the core model. Each is false in general, which is shown by an explicit counterexample, and each has a true core, which is proved under conditions stated exactly. The open question is answered: in the core model, information about compatibility always has non-negative value. So does information about willingness that comes from monotone signals, such as the guide's cues. A non-monotone signal about willingness can have strictly negative value under all three models of the sting. Every number in this note is reproduced by \texttt{verify.py}, which also runs randomised checks of each theorem.
\end{abstract}

\section*{Summary}

\begin{center}
\small
\begin{tabular}{@{}p{2.6cm}p{3.1cm}p{9.0cm}@{}}
\toprule
Claim & Verdict & What is proved \\
\midrule
Conjecture 1 (threshold rules) & False in general (Prop.~\ref{prop:c1-counter}) & Threshold rules hold under a single condition (M) (Theorem~\ref{thm:c1}). The condition is automatic when the after-date choice does not depend on what the date reveals. With two types it reads $v_j(C_H-C_L)\ge\kt\,|\psi(C_H)-\psi(C_L)|$ (Corollary~\ref{cor:c1}). \\[2pt]
Conjecture 2 (outside option) & True for Phase~C; false on the day in general & Raising $U$ shrinks the set of beliefs at which the agent asks for another date. On the day, raising $U$ makes continuing \emph{less} attractive at beliefs where the agent would cut contact after the date, and \emph{more} attractive where they would ask for another date (Lemma~\ref{lem:pivot}, Theorem~\ref{thm:c2}). \\[2pt]
Conjecture 3 (relationship value) & First part: true for Phase~C, false on the day in general. Second part: false & Raising $R$ has the mirror-image effect of raising $U$ (Theorem~\ref{thm:c3}). The investment value of a step peaks where the after-date choice is closest to indifference, not where beliefs are least settled. It vanishes as beliefs settle (Proposition~\ref{prop:inv}). \\[2pt]
Conjecture 4 (tight budgets) & Premise rarely holds; conclusion partly true & Concavity gives the shadow-value claim by definition, but entry prices make $V$ non-concave in $B$. Without any concavity, thresholds rise as the budget falls in the core model. Under concavity, spend falls with the budget. ``Cheaper opening activities'' is false (Proposition~\ref{prop:c4}). \\[2pt]
Conjecture 5 (spending to learn) & First and last parts true; ``peaks under uncertainty'' false & Optimal spend is never below the Proposition~1 level, under all three sting models. It converges to that level as beliefs settle. The premium need not be largest when beliefs are least settled (Proposition~\ref{prop:c5}). \\[2pt]
Open question, Sec.~4.4 & Answered & EVSI $\ge0$ for every signal about $C$, and for every MLRP signal about $q$, under D1, D2 and KR. A non-MLRP signal about $q$ can have negative EVSI under D1, D2 and KR (Theorem~\ref{thm:voi}, Examples~\ref{ex:neg} and~\ref{ex:negD2}). \\
\bottomrule
\end{tabular}
\end{center}

\section{Setting}\label{sec:setting}

Notation is that of the proposal. Throughout, $b$ denotes the agent's belief, $\E_b$ expectation under it, $p=\varphi_j\E_b[q]$ the probability of a yes at step $j$, and
\[
\pp(p)\;=\;p-\kappa_m(p),\qquad m\in\{\mathrm{D1},\mathrm{D2},\mathrm{KR}\},
\]
the \emph{effective continuation weight}: $\kappa_{\mathrm{D1}}=\eta(1-p)$, $\kappa_{\mathrm{D2}}=\eta p(1-p)$, $\kappa_{\mathrm{KR}}=\eta(\Lambda-1)p(1-p)$.

\begin{assumption}[Core model]\label{S:core}
Assumptions A1--A12 hold, including independent priors (A5). The value after the date is truncated after one more date, as in (19)--(21), with $h(c,q_D)=\bar h\,c\,q_D$.
\end{assumption}

\begin{assumption}[Ratings]\label{S:ratings}
After step $j$ the agent observes a rating $y$ whose likelihood $P_j(y\mid c)$ has the monotone likelihood ratio property (MLRP) in $c$, and updates by Bayes' rule. The leading case is $y\sim\mathrm{Bin}(m_j,C)$ with integer $m_j$. Beta beliefs then stay Beta: $\mathrm{Beta}(\alpha,\beta)\mapsto\mathrm{Beta}(\alpha+y,\beta+m_j-y)$.
\end{assumption}

\begin{remark}[The fractional likelihood]
The fractional Binomial (2) is a tempered quasi-likelihood. It does not by itself define a distribution for the rating, so expected posterior moments need not equal prior moments. Under it, $\E[\E_{b'}f(C)\mid b]=\E_bf(C)$ holds for $f(c)=c$ when $\E[\tilde y\mid C]=m_aC$, but not for nonlinear $f$. The results below that use the martingale property for nonlinear functions of $C$ therefore hold exactly with integer $m_a$ and exact Bayes, and approximately with fractional $m_a$. The results that use only monotonicity carry over unchanged.
\end{remark}

\begin{assumption}[No sting without a stake]\label{S:plus}
The expected feeling of an answer is $\kappa_m(p)\,\Delta^{+}$ with $\Delta^{+}=\max\{\Delta,0\}$. The proposal assumes $\Delta>0$ throughout; this extends it to $\Delta\le0$ without turning a refusal into a gain.
\end{assumption}

\begin{assumption}[Willingness to meet again]\label{S:qD}
Where stated, $q_D$ is known, so $P_D=q_D$. This is part of the hypothesis of Conjecture~1. It is used elsewhere only to make the after-date value linear-plus-option in the belief about $C$ (Lemma~\ref{lem:vend}); the same results hold if $q_D$ is uncertain and $\eta=0$ in Phase~C.
\end{assumption}

Write $\kt=\pos{\pp(q_D)}=\pos{q_D-\kappa_m(q_D)}$, and define
\begin{align*}
a(c)&=h(c)\,c\,R+\bigl(1-h(c)\bigr)U,\\
\psi(c)&=\bigl(1-h(c)\bigr)\bigl[\delta_D\bigl(h(c)\,c\,R+(1-h(c))U\bigr)-U\bigr],\qquad h(c)=\bar h\,q_D\,c,
\end{align*}
and the \emph{Phase-C index} $\Psi(b)=\E_b[\psi(C)]$.

\begin{lemma}[After-date value]\label{lem:vend}
Under \ref{S:core}, \ref{S:plus} and \ref{S:qD},
\[
\Vend(b)\;=\;\E_b[a(C)]\;+\;\kt\,\pos{\Psi(b)}.
\]
After a date that did not become a relationship, the agent asks for another date if and only if $\kt>0$ and $\Psi(b)>0$.
\end{lemma}

\begin{proof}
With $q_D$ known, a yes to another date says nothing about $C$, so in (20) $b_{\mathrm{yes}}=b$ and the stake is $\Delta_D=V_D(b_{\mathrm{not}})-U$. If $\Delta_D\ge0$, asking is worth $U+\pp(q_D)\Delta_D$, so $V_C(b_{\mathrm{not}})=U+\kt\Delta_D$. If $\Delta_D<0$, asking is worth $U+q_D\Delta_D<U$, so $V_C=U$. Hence $V_C(b_{\mathrm{not}})=U+\kt\,\Delta_D^{+}$. The event ``no relationship'' has likelihood $1-h(C)$, so $\E_{b_{\mathrm{not}}}[f(C)]\,(1-\E_bh)=\E_b[(1-h(C))f(C)]$ for every $f$. Applying this to (21) gives $(1-\E_bh)(V_D(b_{\mathrm{not}})-U)=\E_b\psi=\Psi(b)$. Since $1-\E_bh>0$, $(1-\E_bh)\Delta_D^{+}=\Psi(b)^{+}$. Substituting into (19):
\[
\Vend(b)=\E_b[hCR]+(1-\E_bh)U+\kt\,\Psi(b)^+=\E_b[a]+\kt\,\Psi(b)^+.\qedhere
\]
\end{proof}

So $\Vend$ is linear in $b$ plus a convex option term. The option has a kink exactly where the after-date choice switches between cutting contact and asking for another date.

\begin{lemma}[A refusal carries no news; effective continuation weight]\label{lem:core}
Under \ref{S:core} and \ref{S:plus}, the loss from a no is $L=0$, so $\ell=0$ and $\Delta=G$, at every step of every date. Writing
\[
G_j(b,B)=\sup_{s\in S_j(B)}\Bigl\{W_j(b;s)-c_j^{\mathrm{time}}+\E\bigl[V_{j+1}(b',B-\sigma(k_j+s))\mid b_{\mathrm{yes}}\bigr]\Bigr\}-\Vend(b)
\]
(with $G_j=-\infty$ if $S_j(B)=\emptyset$), the Bellman equation (8) is equivalent to
\begin{equation}\label{eq:Vj}
V_j(b,B)=\Vend(b)+\pp_j^{\,+}\,G_j(b,B)^{+},\qquad \pp_j=\pp(\varphi_j\E_b q).
\end{equation}
The agent asks if and only if $\pp_j>0$ and $G_j\ge0$, and the optimal venue maximises the term in braces.
\end{lemma}

\begin{proof}
By Lemma~\ref{lem:vend}, or directly from (19)--(21), $\Vend$ depends on $b$ only through its marginal on $(C,q_D)$. By A4 and A5 an answer is informative only about $q$, which is independent of $(C,q_D)$. Hence $\Vend(b_{\mathrm{no}})=\Vend(b_{\mathrm{yes}})=\Vend(b)$, so $L=0$. At spend $s$, let $G(s)$ be the term in braces minus $\Vend(b)$. Asking is worth $\Vend(b)+pG(s)-\kappa(p)G(s)^+$. This is $\Vend+\pp\,G(s)$ when $G(s)\ge0$, and is below $\Vend$ when $G(s)<0$. Maximising over $s$ and comparing with ending gives \eqref{eq:Vj}.
\end{proof}

Since $\pp^{\,+}\in[0,1]$, \eqref{eq:Vj} can be rewritten as
\[
V_j=\max\Bigl\{\Vend(b),\;\pp_j^{\,+}\bigl(G_j+\Vend(b)\bigr)+(1-\pp_j^{\,+})\Vend(b)\Bigr\}.
\]
This is the Bellman equation of an \emph{optimal stopping problem}. After each ask the date continues with probability $\pp^{\,+}$ and pays the step's flow. With probability $1-\pp^{\,+}$ it ends at the current belief about $C$, and when the date ends the agent receives $\Vend$ at the terminal belief. Call this fictitious measure $\Pt$, with expectation $\Et$. Because $\pp$ does not depend on $C$, and ratings are Bayesian, the belief about $C$ is a $\Pt$-martingale. So for every policy $\pi$ and stopping time $\tau\le J+1$,
\begin{equation}\label{eq:mart}
\Et^{\pi}\bigl[\E_{b_\tau}f(C)\bigr]=\E_bf(C)\qquad\text{for every bounded }f.
\end{equation}
This is the tool behind most of what follows.

\begin{remark}[Closed form of Proposition 1]\label{rem:closed}
With $\bar s_j=k_j$ and $x=\E[C]v_j$, the net value of a step at the optimal spend is
\[
W_j=x\ln\!\frac{x}{\sigma\lambda k_j}\quad\text{if }x\ge\sigma\lambda k_j,\qquad W_j=x-\sigma\lambda k_j\quad\text{otherwise.}
\]
Discretionary spend is therefore positive exactly when the step is worth taking at its entry price.
\end{remark}

\section{Conjecture 1: threshold rules}\label{sec:c1}

Conjecture~1 concerns the one-dimensional case, with $q$ and $q_D$ known. Then $p_j=\varphi_jq$ and $\pp_j$ are constants, and the agent asks at step $j$ iff $\pp_j>0$ and $G_j\ge0$ (Lemma~\ref{lem:core}). At each step the belief lies in a totally ordered family $F_j$:
\begin{itemize}[nosep]
\item two types, $b=\pi=P(C=C_H)$; or
\item Beta, $b=\mathrm{Beta}(\alpha,n_j-\alpha)$, with $n_j=n_1+\sum_{i<j}m_i$ fixed by the step.
\end{itemize}
Both families are ordered by the likelihood ratio, and $\E_bC$ is strictly increasing along them. With MLRP ratings:
\begin{enumerate}[label=(\roman*),nosep]
\item the posterior $T_j(b,y)$ is nondecreasing in $b$ and $y$;
\item the predictive distribution of $y$ is increasing in $b$ in the first-order stochastic sense;
\item hence $b\mapsto\E[\phi(T_j(b,y))\mid b]$ is nondecreasing for every nondecreasing $\phi$ on $F_{j+1}$.
\end{enumerate}

\begin{definition}
The \emph{investment value} of step $j$ is
\[
I_j(b)=\E[\Vend(T_j(b,y))\mid b]-\Vend(b)=\kt\Bigl(\E\bigl[\Psi(T_j(b,y))^+\mid b\bigr]-\Psi(b)^+\Bigr)\;\ge0.
\]
The second expression follows from Lemma~\ref{lem:vend} and the martingale property; the sign follows from convexity.
\end{definition}

\begin{theorem}[Threshold rules]\label{thm:c1}
Assume \ref{S:core}--\ref{S:qD} with $q$ known, and assume for every $j$:
\begin{equation*}
\textbf{(M)}\qquad I_j(b_2)-I_j(b_1)\;\ge\;-v_j\bigl(\E_{b_2}C-\E_{b_1}C\bigr)\qquad\text{for all }b_1\le b_2\text{ in }F_j.
\end{equation*}
Then for every step $j$ and budget $B$, $G_j(\cdot,B)$ is nondecreasing on $F_j$. The set of beliefs at which the agent asks is therefore $\{b\in F_j:\E_bC\ge c_j^*(B)\}$, a threshold on $\E[C]$, possibly at $0$ or above $1$.
\end{theorem}

\begin{proof}
Let $O_j=V_j-\Vend=\pp_j^{\,+}G_j^{+}$, and note $O_{J+1}\equiv0$. Suppose $O_{j+1}(\cdot,B')$ is nondecreasing on $F_{j+1}$ for every $B'$; this holds for $j=J$. Since $\Vend$ does not depend on $B$,
\[
G_j(b,B)=\sup_{s\in S_j(B)}\Bigl\{W_j(b;s)-c_j^{\mathrm{time}}+\E\bigl[O_{j+1}(T_j(b,y),B-\sigma(k_j+s))\mid b\bigr]\Bigr\}+I_j(b).
\]
Fix $b_1\le b_2$ and $\varepsilon>0$, and let $s_1$ be $\varepsilon$-optimal at $b_1$. Then:
\begin{itemize}[nosep]
\item $W_j(b_2;s_1)-W_j(b_1;s_1)=v_j\bigl(1+g(s_1)\bigr)(\E_{b_2}C-\E_{b_1}C)\ge v_j(\E_{b_2}C-\E_{b_1}C)$, because $g\ge0$;
\item the $O_{j+1}$ term is nondecreasing in $b$, by (iii) and the induction hypothesis;
\item (M) bounds the fall in $I_j$.
\end{itemize}
Hence $G_j(b_2,B)-G_j(b_1,B)\ge-\varepsilon$ for every $\varepsilon>0$. So $G_j(\cdot,B)$ is nondecreasing, and then so is $O_j(\cdot,B)=\pp_j^{\,+}G_j^+$. Induction completes the proof.
\end{proof}

Condition (M) asks the consumption value of the step to rise with $\E[C]$ at least as fast as the step's investment value can fall. The investment value is the value of what the step reveals for the after-date choice. The next corollary gives conditions on primitives.

\begin{corollary}\label{cor:c1}
Condition (M), and so the conclusion of Theorem~\ref{thm:c1}, holds in each of the following cases.
\begin{enumerate}[label=(\alph*)]
\item \emph{The after-date choice does not depend on what the date reveals:} $\kt=0$, or $\psi$ has constant sign on the support of beliefs. Then $I_j\equiv0$.
\item \emph{Two types:} $v_j(C_H-C_L)\ge\kt\,\lvert\psi(C_H)-\psi(C_L)\rvert$ for every $j$. This does not depend on the rating technology.
\item \emph{Beta beliefs:} $\Psi$ is nondecreasing on every $F_j$, and $\kt L_j\le v_j$, where
\[
L_j=\sup_{b_1<b_2\in F_j}\frac{\Psi(b_2)-\Psi(b_1)}{\E_{b_2}C-\E_{b_1}C}.
\]
If $\psi$ is nondecreasing on $[0,1]$, both parts hold whenever $\kt\sup_c\psi'(c)\le\min_jv_j$.
\end{enumerate}
\end{corollary}

\begin{proof}
(a) If $\psi\ge0$ on the support, $\Psi\ge0$ at every posterior, so $I_j=\kt(\E[\Psi(b')]-\Psi(b))=0$ by \eqref{eq:mart}. The case $\psi\le0$ is the same, and $\kt=0$ is immediate.

(b) Write $x_0=\psi(C_L)$, $x_1=\psi(C_H)$ and $P^H_y=P_j(y\mid C_H)$, $P^L_y=P_j(y\mid C_L)$. Then
\[
I_j(\pi)/\kt=\sum_y\bigl(\pi P_y^Hx_1+(1-\pi)P_y^Lx_0\bigr)^+-\bigl(x_0+\pi(x_1-x_0)\bigr)^+ .
\]
This is continuous and piecewise linear in $\pi$. Where it is differentiable, its slope is
\[
\sum_{y\in A(\pi)}\bigl(P_y^Hx_1-P_y^Lx_0\bigr)-\mathbf 1\{\Psi(\pi)>0\}(x_1-x_0),
\]
where $A(\pi)$ is the set of $y$ whose term is positive.
\begin{itemize}[nosep]
\item If $x_0<0<x_1$, every summand is positive, so the slope is at least $-(x_1-x_0)$.
\item If $x_1<0<x_0$, every summand is negative, so the sum is at least $\sum_y(P_y^Hx_1-P_y^Lx_0)=x_1-x_0$, and the indicator term is non-negative.
\item If $x_0$ and $x_1$ have the same sign, (a) applies.
\end{itemize}
So the slope of $I_j$ in $\pi$ is at least $-\kt|x_1-x_0|$. Since $\E_\pi C=C_L+\pi(C_H-C_L)$, (M) holds under the stated inequality.

(c) For $b_1\le b_2$ in $F_j$, write
\[
\frac{I_j(b_2)-I_j(b_1)}{\kt}=\Bigl[\E\bigl(\Psi^+(b')\mid b_2\bigr)-\E\bigl(\Psi^+(b')\mid b_1\bigr)\Bigr]-\Bigl[\Psi^+(b_2)-\Psi^+(b_1)\Bigr].
\]
The first bracket is non-negative by (iii), because $\Psi^+$ is nondecreasing on $F_{j+1}$. The second is at most $(\Psi(b_2)-\Psi(b_1))^+\le L_j(\E_{b_2}C-\E_{b_1}C)$. If $\psi$ is nondecreasing, first-order dominance gives
\[
\Psi(b_2)-\Psi(b_1)=\int\psi'(c)\,\bigl(F_{b_1}(c)-F_{b_2}(c)\bigr)\,dc\le\sup\psi'\cdot\bigl(\E_{b_2}C-\E_{b_1}C\bigr),
\]
and $\Psi$ is nondecreasing.
\end{proof}

\begin{proposition}[Conjecture 1 fails without (M)]\label{prop:c1-counter}
Take two types $C\in\{0.2,0.8\}$ and one step left, with $v=50$, $k=30$, $c^{\mathrm{time}}=0$, $\sigma=\lambda=1$ and one Bernoulli rating ($m=1$). Let $q=0.9$, D1 with $\eta=0.2$, $q_D=0.8$, $\bar h=0.5$, $\delta_D=0.9$, $R=6000$, $U=1000$. Then the agent asks at the last step on $(0.071,\,0.232)\cup(0.667,\,1]$ and ends on the rest of $[0,1]$. The set of asking beliefs is not an interval, and no threshold rule is optimal.
\end{proposition}

\begin{proof}
The primitives give:
\begin{itemize}[nosep]
\item $\psi(0.2)=-78.75$ and $\psi(0.8)=676.19$, so $\Psi(\pi)=0$ at $\pi_k=0.1043$;
\item $\kt=0.8-0.2\cdot0.2=0.76$ and $\pp=0.9-0.2\cdot0.1>0$;
\item $s^*=\max(50\,\E C-30,0)$, which is $0$ for $\E C\le0.6$, so $W(\pi)=50\,\E_\pi C-30$ there.
\end{itemize}
Consider four regions.
\begin{itemize}
\item For $\pi\ge0.318$, the posterior after a low rating, $\pi/(\pi+4(1-\pi))$, is above $\pi_k$. So $\Psi^+$ is linear on the support of the posteriors, and $I(\pi)=0$. Then $G=W<0$ for $\E_\pi C<0.6$, that is, for $\pi<0.667$; and $G\ge0$ above.
\item At $\pi=\pi_k$, only a high rating leads to asking for another date. So
\[
I(\pi_k)=\kt\bigl(\pi_k\cdot0.8\cdot676.19+(1-\pi_k)\cdot0.2\cdot(-78.75)\bigr)=32.17,
\]
while $W(\pi_k)=-16.87$, giving $G(\pi_k)=15.29>0$.
\item At $\pi=0.05$, $G=-9.32<0$.
\end{itemize}
$G$ therefore changes sign three times. \texttt{verify.py} computes $G$ on a fine grid and confirms the switch points.
\end{proof}

The mechanism is the after-date option. Near $\pi_k$ the agent is close to indifferent between cutting contact and asking for a second date. The last step's rating settles that choice, so the step is worth taking for what it reveals, even though it is not worth its price as an evening out. At somewhat higher beliefs the second date is already decided, the information is worthless, and the agent ends. Here (M) fails badly: $v(C_H-C_L)=30$, while $\kt|\psi(C_H)-\psi(C_L)|=574$.

\begin{remark}
When (M) fails, the asking set at each step is still a finite union of intervals in $\E[C]$. A card can record it as such, with extra columns near the after-date indifference point.
\end{remark}

\section{Conjectures 2 and 3: the outside option and the value of a relationship}\label{sec:c23}

$U$ and $R$ enter the core model only through $\Vend$, that is, through $a$ and $\psi$. One lemma covers both.

\begin{lemma}[Pivotality]\label{lem:pivot}
Assume \ref{S:core}--\ref{S:qD}. Beliefs about $q$ and the budget are arbitrary. Let a parameter change from $\theta$ to $\theta'$ change $a$ arbitrarily and raise $\psi$ pointwise, $\psi_{\theta'}\ge\psi_\theta$ on $[0,1]$, leaving everything else unchanged. Let
\[
D_j=V^{\mathrm{ask}}_j-\Vend
\]
be the value of asking minus the value of ending at step $j$; when $\pp_j>0$, $D_j=\pp_jG_j$.
\begin{enumerate}[label=(\roman*)]
\item If $\Psi_\theta(b)>0$, so that the agent would ask for another date, then $D_j(b;\theta')\le D_j(b;\theta)$.
\item If $\Psi_{\theta'}(b)\le0$, so that the agent would cut contact, then $D_j(b;\theta')\ge D_j(b;\theta)$.
\end{enumerate}
\end{lemma}

\begin{proof}
By Lemma~\ref{lem:core}, the value of any policy $\pi$ that asks now is
\[
\Phi(\pi;\theta)=\Et^{\pi}\Bigl[\textstyle\sum\text{flows}\Bigr]+\Et^{\pi}\bigl[\Vend(b_\tau;\theta)\bigr].
\]
Neither the flows nor $\Pt^{\pi}$ depend on $\theta$; only the terminal payoff does. Let $\pi$ and $\pi'$ be $\varepsilon$-optimal asking policies at $\theta$ and $\theta'$, and write $\delta V(x)=\Vend(x;\theta')-\Vend(x;\theta)$. Then
\[
\Et^{\pi}[\delta V(b_\tau)]-\varepsilon\;\le\;V_j^{\mathrm{ask}}(\theta')-V_j^{\mathrm{ask}}(\theta)\;\le\;\Et^{\pi'}[\delta V(b_\tau)]+\varepsilon.
\]
By Lemma~\ref{lem:vend}, $\delta V(x)=\E_x[a_{\theta'}-a_\theta]+\kt\,\phi(x)$, where
\[
\phi(x)=\Psi_{\theta'}(x)^+-\Psi_\theta(x)^+,\qquad 0\le\phi(x)\le\E_x\zeta,\qquad \zeta=\psi_{\theta'}-\psi_\theta\ge0.
\]
The first term of $\delta V$ is linear in $x$, so by \eqref{eq:mart} it cancels against the same term in $\delta V(b)$.

(i) Here $\Psi_{\theta'}(b)\ge\Psi_\theta(b)>0$, so $\phi(b)=\E_b\zeta$. Then
\[
D_j(\theta')-D_j(\theta)\le\kt\bigl(\Et^{\pi'}[\E_{b_\tau}\zeta]-\E_b\zeta\bigr)+\varepsilon=\varepsilon,
\]
using \eqref{eq:mart} again.

(ii) Here $\Psi_\theta(b)\le\Psi_{\theta'}(b)\le0$, so $\phi(b)=0$. Then
\[
D_j(\theta')-D_j(\theta)\ge\kt\,\Et^{\pi}[\phi(b_\tau)]-\varepsilon\ge-\varepsilon.
\]
Let $\varepsilon\to0$.
\end{proof}

The intuition is that the direct value of $\Vend$ is linear in beliefs, so it is the same, in expectation, whether the agent continues or ends. A parameter therefore changes the value of continuing only through the information value of the date for the after-date choice. That value rises when the parameter moves the after-date choice towards indifference, and falls when it moves the choice away.

\begin{theorem}[Conjecture 2]\label{thm:c2}
Assume \ref{S:core} and \ref{S:plus}, and let $U<U'$.
\begin{enumerate}[label=(\alph*)]
\item \emph{Phase C (true).} The set of beliefs at which the agent asks for another date shrinks. The agent asks iff $\pp(P_D)>0$ and $V_D(b_{\mathrm{yes}})>U$, and
\[
V_D-U=\delta_D\E[hC]R-\bigl(1-\delta_D(1-\E h)\bigr)U
\]
is strictly decreasing in $U$. This part does not need \ref{S:qD}.
\item \emph{On the day (true only where the agent would cut contact).} Assume \ref{S:qD} and let $A_j(U)$ be the set of beliefs at which the agent asks at step $j$.
\begin{itemize}[nosep]
\item On $\{b:\Psi_U(b)\le0\}$, where the agent would cut contact after the date: $A_j(U')\subseteq A_j(U)$. Raising $U$ makes continuing harder, as conjectured.
\item On $\{b:\Psi_{U'}(b)>0\}$, where the agent would ask for another date even at $U'$: $A_j(U)\subseteq A_j(U')$. Raising $U$ makes continuing \emph{easier}, contrary to the conjecture.
\end{itemize}
\item \emph{Spend.} While the budget does not bind, discretionary spend does not depend on $U$, by Proposition~1.
\end{enumerate}
\end{theorem}

\begin{proof}
(b) We have $\partial\psi/\partial U=-(1-h)\bigl(1-\delta_D(1-h)\bigr)\le0$, so $\psi_U\ge\psi_{U'}$. Apply Lemma~\ref{lem:pivot} with $\theta=U'$ and $\theta'=U$, recalling that the agent asks iff $\pp_j>0$ and $D_j\ge0$, and that $\pp_j$ does not depend on $U$.
\end{proof}

\begin{example}[Conjecture 2 fails on the day, inside a threshold model]\label{ex:c2}
Use the model of Proposition~\ref{prop:c1-counter} with $R=600$, $U=100$ and a last step with $v=100$ and $k=32$. Corollary~\ref{cor:c1}(b) holds, since $60\ge57.4$, so the asking set is a threshold.
\begin{itemize}[nosep]
\item The threshold is $\pi^*=0.1605$, where $\Psi=4.24>0$.
\item Raising $U$ by $5$ \emph{lowers} the threshold to $0.1508$.
\item With $k=25$ instead, the threshold is $0.0606$, where $\Psi=-3.30<0$. There raising $U$ by $5$ raises it to $0.0622$, as conjectured.
\end{itemize}
\end{example}

\begin{theorem}[Conjecture 3, first part]\label{thm:c3}
Assume \ref{S:core} and \ref{S:plus}, and let $R<R'$.
\begin{enumerate}[label=(\alph*)]
\item \emph{Phase C (true).} The set of beliefs at which the agent asks for another date grows, because $V_D-U$ is increasing in $R$.
\item \emph{On the day.} Assume \ref{S:qD}.
\begin{itemize}[nosep]
\item On $\{b:\Psi_R(b)>0\}$: $A_j(R')\subseteq A_j(R)$. A more valuable relationship makes continuing \emph{harder}, contrary to the conjecture.
\item On $\{b:\Psi_{R'}(b)\le0\}$: $A_j(R)\subseteq A_j(R')$, as conjectured.
\end{itemize}
In Example~\ref{ex:c2}, raising $R$ by $30$ raises the threshold from $0.1605$ to $0.1677$. With $k=25$ it lowers the threshold from $0.0606$ to $0.0584$.
\end{enumerate}
\end{theorem}

\begin{proof}
$\partial\psi/\partial R=(1-h)\,\delta_D\,h\,c\ge0$. Apply Lemma~\ref{lem:pivot} with $\theta=R$ and $\theta'=R'$.
\end{proof}

\begin{remark}[Why $R$ does not simply favour escalation]
In the core model the chance of a relationship, $h(C,q_D)$, does not depend on how far up the ladder the date went. The direct value of a relationship, $\E[hC]R$, is therefore a martingale term that continuing cannot raise. If the intended reading is that escalating raises the chance of a relationship, $h$ must depend on the step reached. Conjecture~3 would then need a separate argument.
\end{remark}

\begin{proposition}[Conjecture 3, second part: investment value]\label{prop:inv}
Assume \ref{S:core}--\ref{S:qD}.
\begin{enumerate}[label=(\roman*)]
\item $I_j\ge0$.
\item \emph{Two types}, with $x_0=\psi(C_L)<0<x_1=\psi(C_H)$: $I_j$ is nondecreasing on $[0,\pi_k]$ and nonincreasing on $[\pi_k,1]$, where $\pi_k=-x_0/(x_1-x_0)$ is the after-date indifference point. Moreover $I_j(\pi)=0$ whenever every posterior lies on the same side of $\pi_k$ as $\pi$. If $x_0,x_1$ have the same sign, $I_j\equiv0$.
\item \emph{Settling:} for Beta beliefs, writing $\psi(c)=\sum_k\psi_kc^k$, which has degree at most $3$,
\[
I_j\bigl(\mathrm{Beta}(\alpha,n-\alpha)\bigr)\le\kt\,\frac{m_j}{n}\sum_{k\le3}k\lvert\psi_k\rvert\;\longrightarrow\;0\qquad(n\to\infty).
\]
\item \emph{``Largest when least settled'' is false.} In the model of Proposition~\ref{prop:c1-counter}, $I(1/2)=0$ but $I(\pi_k)=32.2$. The consumption value depends on beliefs only through $\E[C]$, so the investment share of the step's value is zero at maximal uncertainty, $\pi=1/2$. It is positive at $\pi_k\approx0.10$, where the agent is fairly sure the match is incompatible. Nor is $I_j$ monotone in precision at a fixed mean: with Beta beliefs of mean $0.3$, $I=4.14$ at $n=8$ but $4.72$ at $n=16$.
\end{enumerate}
\end{proposition}

\begin{proof}
(ii) On $\{\Psi\le0\}$, $I/\kt=\sum_y(\cdot)^+$ with every summand nondecreasing, because its slope $P_y^Hx_1-P_y^Lx_0$ is positive. On $\{\Psi\ge0\}$, $I/\kt=\sum_y(\cdot)^+-\sum_y(\cdot)=\sum_y(\cdot)^-$, with every summand nonincreasing. If all posteriors lie on one side of $\pi_k$, $(\cdot)^+$ is linear on them and \eqref{eq:mart} gives $I=0$.

(iii) Since $x\mapsto x^+$ is $1$-Lipschitz, $I_j\le\kt\,\E|\Psi(b')-\Psi(b)|$. For $b'=\mathrm{Beta}(\alpha+y,n-\alpha+m-y)$, each factor of $\E_{b'}C^k=\prod_{i<k}\frac{\alpha+y+i}{n+m+i}$ differs from the corresponding factor of $\E_bC^k$ by at most $m/(n+m+i)\le m/n$. All factors lie in $[0,1]$, so $|\E_{b'}C^k-\E_bC^k|\le km/n$.

(iv) From $\pi=1/2$ the posteriors are $0.2$ and $0.8$, both above $\pi_k$, so (ii) gives $I=0$. The other values are computed in \texttt{verify.py}.
\end{proof}

This is the classical point that the value of information is not monotone in prior uncertainty \cite{gould,hilton}. The true statement is that the investment share is largest near the after-date indifference point, and that it vanishes as beliefs settle.

\section{Conjecture 4: tight budgets}\label{sec:c4}

\begin{proposition}\label{prop:c4}
Assume \ref{S:core} and \ref{S:plus}.
\begin{enumerate}[label=(\roman*)]
\item \emph{Shadow value (true, by definition).} If $V(b,\cdot)$ is concave, $\mu(B)=[V(B)-V(B-\varepsilon)]/\varepsilon$ is nonincreasing in $B$.
\item \emph{The premise usually fails.} Suppose step $j$ has entry price $k_j>0$, $\pp_j>0$, and $G_j(b,\sigma k_j)>0$. Then $V_j(b,\cdot)$ jumps up at $B=\sigma k_j$ and is not concave on any interval containing $\sigma k_j$ in its interior.
\item \emph{Thresholds (true in the core model, with no concavity needed).} $G_j(b,\cdot)$ is nondecreasing in $B$ for every $b$ and $j$. So the set of beliefs at which the agent asks shrinks as the budget falls.
\item \emph{Spend (true under concavity).} If $V_{j+1}(b',\cdot)$ is concave for every $b'$, every selection of the optimal spend $s_j^*(b,B)$ can be taken nondecreasing in $B$. With less budget, the agent spends less, as if money were more expensive.
\item \emph{Cheaper opening activities (false).} The cost of the optimal opening activity need not rise with the budget, nor fall with the price of money. See Examples~\ref{ex:openB} and~\ref{ex:openL}.
\item \emph{Beyond the core.} Once the value after the date uses the budget, as in M2+ where another date must be paid for, thresholds can \emph{fall} as the budget falls. See Example~\ref{ex:m2plus}.
\end{enumerate}
\end{proposition}

\begin{proof}
(i) Slopes of chords of a concave function are nonincreasing.

(ii) For $B<\sigma k_j$ the step is unaffordable and $V_j=\Vend$. At $B=\sigma k_j$, $V_j=\Vend+\pp_jG_j>\Vend$. A nondecreasing function with an upward jump is not concave across it.

(iii) By induction, $V_{j+1}(b',\cdot)$ is nondecreasing: $V_{J+1}=\Vend$ does not depend on $B$, and \eqref{eq:Vj} preserves monotonicity. The feasible set $S_j(B)$ grows with $B$, and $\Vend$ does not depend on $B$.

(iv) For $s<s'$, $F(s',B)-F(s,B)$ consists of a part independent of $B$ plus
\[
\E\bigl[V_{j+1}(x)-V_{j+1}(x+d)\bigr],\qquad x=B-\sigma(k_j+s'),\quad d=\sigma(s'-s).
\]
This is nondecreasing in $B$ by concavity. So $F$ has increasing differences in $(s,B)$, $S_j(B)=[0,B/\sigma-k_j]$ is ascending, and Topkis' theorem applies.
\end{proof}

\begin{example}[Opening activity and the budget]\label{ex:openB}
Let $C$ be known, $q=1$, $\eta=0$ and $\sigma=\lambda=1$. Coffee has $x=\E[C]v=10$ and $k=8$; dinner has $x=60$ and $k=40$. Spend is optimal subject to the budget. A coffee opening leads to coffee then dinner; a dinner opening leads to dinner only.
\begin{itemize}[nosep]
\item $B=45$: coffee-first is worth $2.23$, because dinner is no longer affordable after coffee. Dinner-first is worth $22.07$. The agent opens with dinner.
\item $B=80$: coffee-first is worth $26.56$ and dinner-first $24.33$. The agent opens with coffee.
\end{itemize}
The richer agent chooses the cheaper opening. This $V(B)$ is not concave, as (ii) predicts.
\end{example}

\begin{example}[Opening activity and the price of money]\label{ex:openL}
Read ``as if money were more expensive'' as a rise in $\lambda$. Let coffee have $x=10.8$ and $k=8$, dinner $x=63$ and $k=40$, with $\sigma=1$. Dinner is worth taking in either ladder for all $\lambda$ near $\lambda_0=1.35$. Coffee-first minus dinner-first equals the coffee's own net value $W_c(\lambda)$ (Remark~\ref{rem:closed}). This is $0.408$ at $\lambda=1.30$ and $-0.4$ at $\lambda=1.40$. Dearer money switches the opening from coffee to dinner.

By the envelope theorem, $\partial V^{a}/\partial\lambda=-\sigma\,\E^{a}[\text{outlay}]$. A higher $\lambda$ therefore favours the opening with the lower \emph{expected total outlay}, which need not be the cheaper opening.
\end{example}

\begin{remark}
Under the conjecture's own premise, the overall value $V=\max(V^{\mathrm{coffee}},V^{\mathrm{dinner}})$ is concave. Where the optimal opening switches transversally, a maximum of two differentiable functions has a convex kink. So the premise allows switches only at points of tangency, and cannot by itself deliver a systematic move to cheaper openings.
\end{remark}

\begin{example}[M2+: thresholds can fall with the budget]\label{ex:m2plus}
At the last step the net value is $w=10$ and the entry price $k=20$, with $\sigma=1$. Suppose the after-date option is worth $30$ only if at least $40$ of budget remains for another date. Then
\[
G_J(B)=w+\Vend(B-20)-\Vend(B).
\]
This is $10$ at $B=30$, where a second date is unaffordable anyway; $-20$ at $B=50$, where continuing would cost the second date; and $10$ at $B=70$. The agent asks at $30$, ends at $50$, and asks at $70$.
\end{example}

\section{Conjecture 5: spending to learn}\label{sec:c5}

Relax A7 as in Appendix~A.7: the rating after step $j$ comes from a family $\{P_m\}$ that is increasing in $m$ in Blackwell's order, with $m=m(s)$ nondecreasing in spend. Integer Binomial $m$ qualifies, as does Gaussian precision. Write
\[
f_\mu(s)=\mu v_jg(s)-\sigma\lambda s .
\]
Proposition~1's spend $s^0(\mu)$ is the unique maximiser of $f_\mu$ on $[0,\infty)$, which is strictly concave when $\mu v_j>0$.

\begin{proposition}\label{prop:c5}
Assume \ref{S:core}--\ref{S:qD}, a slack budget and $\pp_j>0$. Beliefs about $q$ are arbitrary, and the sting model is any of D1, D2, KR.
\begin{enumerate}[label=(\roman*)]
\item \emph{Optimal spend is never below Proposition~1's spend.} Every optimal spend satisfies $s_j^*(b)\ge s^0(\E_bC)$, under all three sting models.
\item \emph{It falls back as beliefs settle.} Suppose $m(\cdot)$ takes finitely many values, and beliefs settle: $b_n\to\delta_{\bar c}$ weakly, for example Beta beliefs with $n\to\infty$ and mean $\to\bar c$. Then $s_j^*(b_n)\to s^0(\bar c)$.
\item \emph{``Peaks under uncertainty'' is false.} Take two types and the after-date model of Proposition~\ref{prop:c1-counter}, at the last step, with spend buying a binary signal of accuracy
\[
a(s)=\tfrac12+0.35(1-e^{-s/20})<1-\pi_k .
\]
At $\pi=1/2$, where uncertainty is greatest, the information premium is exactly zero: $s^*=s^0=0$. At $\pi_k=0.104$, $s^*=18.85$ while $s^0=0$.
\end{enumerate}
\end{proposition}

\begin{proof}
(i) By Lemma~\ref{lem:core}, the optimal venue maximises $\Phi(s)=f_\mu(s)+e(s)$, with
\[
e(s)=\E\bigl[V_{j+1}(b'_{m(s)})\mid b\bigr].
\]
$V_{j+1}$ is convex in the belief about $C$ (Theorem~\ref{thm:voi}(a)), so Blackwell's theorem makes $e$ nondecreasing. If $s<s^0$, then $f_\mu(s)<f_\mu(s^0)$ and $e(s)\le e(s^0)$, so $s$ is not optimal.

(ii) We first show $V_{j+1}$ is continuous in $b$ for the weak topology on beliefs about $C$, including degenerate beliefs. Each ingredient is: $\Vend(b)=\E_b a+\kt\Psi(b)^+$; $W_j$, which depends on $\E_bC$; and the preposterior map $b\mapsto\sum_yP(y\mid b)V(b'_y)$. In the last of these, $P(y\mid b)=\E_b[P(y\mid C)]$ is continuous, the unnormalised posterior is weakly continuous, and $|P(y\mid b)V(b'_y)|\le P(y\mid b)\sup|V|$ near beliefs with $P(y\mid b)=0$.

Hence $\mathrm{EVSI}_m(b)=e_m(b)-V_{j+1}(b)$ is continuous, non-negative, and zero at degenerate beliefs. So
\[
\varepsilon_n=\max_m\mathrm{EVSI}_m(b_n)\to0 .
\]
Maximisers of $\Phi_n=f_{\mu_n}+e$ lie in a fixed compact set, because $f_{\mu_n}\to-\infty$ linearly while $e$ varies by at most $\varepsilon_n$. On that set, $f_{\mu_n}\to f_{\bar c}$ uniformly. Since $f_{\bar c}$ is strictly concave with unique maximiser $s^0(\bar c)$, the standard argument applies: for $|s-s^0|\ge\epsilon$, $f_{\bar c}(s)\le f_{\bar c}(s^0)-\eta_\epsilon$. Once $2\varepsilon_n+2\sup|f_{\mu_n}-f_{\bar c}|<\eta_\epsilon$, every maximiser of $\Phi_n$ is within $\epsilon$ of $s^0$.

(iii) From $\pi=1/2$ the low posterior is $1-a(s)>\pi_k$, so both posteriors lie above the kink and $\mathrm{EVSI}\equiv0$ for every $s$. At $\pi_k$, $\mathrm{EVSI}(s)=\kt(2a(s)-1)\,\pi_k\,\psi(C_H)$ is strictly increasing in $s$. The optimal spends are computed in \texttt{verify.py}.
\end{proof}

\begin{remark}
The proposal's caveat that ``with psychological costs of refusal EVSI is not guaranteed to be non-negative'' does not bite here. The information that spend buys in A.7 is about compatibility, and Theorem~\ref{thm:voi}(a) shows that such information is always weakly valuable in the core model, under all three sting models.
\end{remark}

\section{The value of information (Section 4.4)}\label{sec:voi}

\begin{definition}
Fix a belief $b_C$ about $C$. A function $V$ of beliefs about $q$ is in the class $\mathcal C$ if it is nondecreasing in the likelihood-ratio order, and convex along every \emph{MLR segment}, that is, along $t\mapsto(1-t)b_1+tb_2$ for $b_1\le_{\mathrm{LR}}b_2$.
\end{definition}

\begin{theorem}[Sign of the value of information]\label{thm:voi}
Assume \ref{S:core}--\ref{S:qD}.
\begin{enumerate}[label=(\alph*)]
\item \emph{Information about compatibility.} For every $j$, every $B$ and every belief about $q$, $V_j$ is convex in the belief about $C$. So every signal about $C$ has non-negative EVSI under D1, D2 and KR. This includes ratings and the opening activity's signal in Definition~2.
\item \emph{Information about willingness.} Suppose the sting satisfies
\begin{equation*}
\textbf{(K)}\qquad
\begin{gathered}
f(p)=p-\kappa(p)\text{ is convex and twice differentiable,}\\
\text{and }\kappa(p)/p\text{ is nonincreasing on }(0,1].
\end{gathered}
\end{equation*}
D1, D2 and KR all satisfy (K). Suppose also that the cues observed during the date are MLRP signals about $q$. Then every MLRP signal about $q$ has non-negative EVSI at every decision point of the date.
\item \emph{Non-monotone information about willingness can hurt.} Under D1, D2 and KR, some signals about $q$ that are not MLR-ordered have strictly negative EVSI (Examples~\ref{ex:neg} and~\ref{ex:negD2}).
\end{enumerate}
\end{theorem}

\begin{proof}[Proof of (a)]
We argue by backward induction using Lemma~\ref{lem:core}, writing
\[
V_j=\max\Bigl\{\Vend,\;\pp_j^{\,+}\bigl(\sup_s[W_j(b;s)-c_j+\E(V_{j+1}(b')\mid b)]\bigr)+(1-\pp_j^{\,+})\Vend\Bigr\}.
\]
Each ingredient is convex in $b_C$:
\begin{itemize}[nosep]
\item $\Vend$ is convex (Lemma~\ref{lem:vend});
\item $W_j(b;s)$ is affine in $b_C$;
\item $b_C\mapsto\sum_yP(y\mid b_C)V_{j+1}(b'_y)$ is a sum of perspectives of the convex function $V_{j+1}$, hence convex;
\item a supremum of convex functions is convex;
\item $\pp_j^{\,+}\in[0,1]$ does not depend on $b_C$, so the middle term is a convex combination of convex functions.
\end{itemize}
A maximum of convex functions is convex. Convexity of the value in beliefs is equivalent to non-negative EVSI for every signal.
\end{proof}

\begin{proof}[Proof of (b)]
Three facts about MLR segments are standard. Let $b_t=(1-t)b_1+tb_2$ with $b_1\le_{\mathrm{LR}}b_2$.
\begin{itemize}[nosep]
\item $b_t$ is increasing in $t$ in the likelihood-ratio order.
\item For any likelihood $\ell(z\mid q)$, the posterior $b_{t,z}$ lies on the MLR segment from $b_{1,z}$ to $b_{2,z}$, with weight $tP(z\mid b_2)/P(z\mid b_t)$ on $b_{2,z}$.
\item In particular, after a yes ($\ell=\varphi q$), the posterior lies on the segment from $b_{1,\mathrm{yes}}$ to $b_{2,\mathrm{yes}}$ at weight $s(t)=t\bar p_2/\bar p(t)$, where $\bar p(t)=\E_{b_t}q$ is affine and nondecreasing in $t$.
\end{itemize}

\emph{Step 1: expectations over MLRP signals preserve $\mathcal C$.} Let $V\in\mathcal C$ and let $z$ be an MLRP signal. Along a segment,
\[
\E[V(b_{t,z})]=\sum_z\bigl(A_z(t)+B_z(t)\bigr)\,\phi_z\!\Bigl(\frac{B_z(t)}{A_z(t)+B_z(t)}\Bigr),
\]
with $A_z(t)=(1-t)P(z\mid b_1)$ and $B_z(t)=tP(z\mid b_2)$ affine, and $\phi_z$ convex. This is a perspective, hence convex in $t$. Monotonicity follows from MLRP. Expectations over ratings (independent of $q$) and suprema over spend also preserve $\mathcal C$.

\emph{Step 2: the ask step preserves $\mathcal C$.} Suppose $X\in\mathcal C$, and recall from Lemma~\ref{lem:core} that $V_j-\Vend=f(p)^+X(b_{\mathrm{yes}})^+$.

Monotonicity: by (K), $f=p\cdot(f/p)$ is nondecreasing where positive; $p$ is nondecreasing in the order; and $b\mapsto b_{\mathrm{yes}}$ is monotone.

Convexity along a segment: let $g(t)=f(p(t))\,\phi(s(t))$, where $\phi(s)=X$ along the yes-segment is convex and nondecreasing. Assume $\phi\in C^2$ for now. Since $s'=\bar p_1\bar p_2/\bar p^2\ge0$ and $s''=-2\bar p's'/\bar p$,
\[
g''=f''\,p'^2\,\phi+f\,\phi''\,s'^2+2\,p'\phi's'\Bigl(f'-\frac fp\Bigr).
\]
On the set where $f>0$ and $\phi>0$, every term is non-negative:
\begin{itemize}[nosep]
\item $f''\ge0$ by (K);
\item $p'\ge0$, because LR order implies ordered means;
\item $\phi'\ge0$ and $\phi''\ge0$;
\item $f'-f/p=p\,(f/p)'\ge0$ by (K).
\end{itemize}
A general convex nondecreasing $\phi$ is a uniform limit of $C^2$ ones, and convexity survives limits. The set where $f>0$ and $\phi>0$ is an interval $[t_0,1]$, and $g$ vanishes at $t_0$ if $t_0>0$. So $f^+\phi^+$ is $0$ on $[0,t_0]$ and convex nondecreasing on $[t_0,1]$, hence convex on $[0,1]$.

Induction on $j$, starting from $V_{J+1}-\Vend\equiv0$, gives $V_j(b_C,\cdot)-\Vend(b_C)\in\mathcal C$.

\emph{Step 3: from convexity on segments to non-negative EVSI.} Let $z\in\{z_1<\dots<z_K\}$ be an MLRP signal. Reveal it in stages: first whether $z=z_K$, then, given $z<z_K$, whether $z=z_{K-1}$, and so on. At each stage MLRP implies $b_{\{z=z_k\}}\ge_{\mathrm{LR}}b_{\{z<z_k\}}$, and the current belief lies on the MLR segment between them. So the expected value weakly rises at each stage. The tower property gives $\E[V(b_z)]\ge V(b)$.
\end{proof}

\begin{example}[Negative EVSI under D1]\label{ex:neg}
Let $C$ be known, normalise $\Vend=0$, and take two steps left with net values $w_1=0$ and $w_2=1$. Let $\varphi_2=1$, D1 with $\eta=1/2$, so $f(p)=1.5p-0.5$. Let $q\in\{0.2,0.6,1\}$ with prior $(0.3,0.5,0.2)$. The signal reveals whether $q$ is extreme, $\{0.2,1\}$, or middling, $\{0.6\}$. It is not MLR-ordered. The value is
\[
V(b)=f(\E_bq)\bigl(w_1+f(\E_b[q\mid\mathrm{yes}])\,w_2\bigr).
\]
\begin{itemize}[nosep]
\item Prior: $\E q=0.56$ and $\E[q\mid\mathrm{yes}]=0.7$, so $V=0.34\times0.55=0.1870$.
\item Posterior ``extreme'', probability $1/2$, $b_1=(0.6,0,0.4)$: $V=0.28\times0.7231=0.2025$.
\item Posterior ``middling'', $b_2=\delta_{0.6}$: $V=0.4\times0.4=0.16$.
\end{itemize}
So $\mathrm{EVSI}=0.1812-0.1870=-0.0058<0$. With $\eta=0$ the EVSI is exactly zero.

The signal moves the chance of a yes now, $\E q$, and the chance of a yes later, $\E[q\mid\mathrm{yes}]$, in opposite directions. That breaks the third term of $g''$ in the proof of (b).
\end{example}

\begin{example}[Negative EVSI under D2 and KR]\label{ex:negD2}
Let $q\in\{0.1,0.8,1\}$ with prior $(\tfrac14,\tfrac12,\tfrac14)$, and use the same ``extreme versus middling'' signal. Take D2 with $\eta=2$, equivalently KR with $\eta(\Lambda-1)=2$, so $f(p)=p(2p-1)$; and take $w_1=-\tfrac12$, $w_2=1$, $\varphi_2=1$.
\begin{itemize}[nosep]
\item Prior: $V=0.2363\times0.0906=0.0214$.
\item Posterior ``extreme'': $V=0.055\times0.2679=0.0147$.
\item Posterior ``middling'': $V=0$, because step 2 is no longer worth asking for and step 1 is not worth taking alone.
\end{itemize}
So $\mathrm{EVSI}=0.0074-0.0214=-0.0140<0$.
\end{example}

\paragraph{Answer to the open question.}
In the core model, the value of information can be negative only for information about willingness that is not ordered by the likelihood ratio. Ratings, and cues of the kind used in Section~5.4, which are MLRP, always have non-negative value. This is consistent with the proposal's numerical search, which found no counterexample.

\section{The conjectural prediction about opening activities}

Prediction~6 of Appendix~A.9 says that more uncertainty before the date favours a more informative opening activity, depending on where beliefs sit relative to the decision threshold. As an unconditional statement this is false, for the same reason as Proposition~\ref{prop:inv}(iv): at maximal uncertainty the EVSI of both openings can be zero. What is true follows from Theorem~\ref{thm:voi}(a), Blackwell's theorem and Proposition~\ref{prop:inv}:
\begin{itemize}[nosep]
\item the informativeness premium of a Blackwell-more-informative opening is never negative;
\item it vanishes as beliefs settle;
\item with two types it is largest near the relevant indifference point.
\end{itemize}

\section{Suggested changes to the proposal}

\begin{enumerate}
\item Replace Conjecture~1 by Theorem~\ref{thm:c1} with Corollary~\ref{cor:c1}, and cite Proposition~\ref{prop:c1-counter}. Where (M) fails, the card may need more than one boundary near the after-date indifference point.
\item Replace Conjectures~2 and~3 (first part) by Lemma~\ref{lem:pivot} and Theorems~\ref{thm:c2}--\ref{thm:c3}. Keep the Phase~C statements, which are true. Note that in the core model $R$ affects escalation only through information, because $h$ does not depend on the step reached. Prediction~2 (market thickness) and Prediction~5 (relationship value) of Appendix~A.9 should be restated accordingly.
\item Replace ``largest when beliefs are least settled'', in Conjectures~3 and~5 and in Prediction~6, by ``largest near the after-date indifference point, and vanishing as beliefs settle''.
\item In Conjecture~4, drop the concavity premise, which fails whenever entry prices bind. State Proposition~\ref{prop:c4}(iii)--(iv) instead, and drop ``cheaper opening activities'' from the conjecture and from Prediction~3.
\item In Section~4.4 and the ``Convexity'' bullet of Section~8, replace the open question with Theorem~\ref{thm:voi}.
\item Technical points:
\begin{itemize}[nosep]
\item state the convention $\kappa\Delta^+$;
\item note that the fractional likelihood is a quasi-likelihood, so results relying on the martingale property are exact only with integer $m_a$;
\item record the closed form in Remark~\ref{rem:closed};
\item record that $\ell=0$ throughout the core model (Lemma~\ref{lem:core}), so the decision to ask on the day reduces to $p\ge p^*_m(0)$ together with $G\ge0$.
\end{itemize}
\end{enumerate}

\section*{Numerical verification}
\texttt{python3 verify.py} reproduces every number above and runs randomised checks of:
\begin{itemize}[nosep]
\item Corollary~\ref{cor:c1}(b)--(c): monotone $G_j$ whenever the condition holds;
\item the sign pattern of Lemma~\ref{lem:pivot} for $U$ and $R$, at several thousand belief points;
\item Proposition~\ref{prop:c4}(iii);
\item Theorem~\ref{thm:voi}(a), as convexity of $V_0$ in beliefs about $C$;
\item Theorem~\ref{thm:voi}(b), as non-negative EVSI of Binomial cues about $q$ under general priors.
\end{itemize}
The output is saved in \texttt{verify\_output.txt}.

\begin{thebibliography}{9}
\bibitem{gould} Gould, J.~P. (1974). Risk, stochastic preference, and the value of information. \emph{Journal of Economic Theory} 8(1), 64--84.
\bibitem{hilton} Hilton, R.~W. (1981). The determinants of information value: synthesizing some general results. \emph{Management Science} 27(1), 57--64.
\bibitem{topkis} Topkis, D.~M. (1978). Minimizing a submodular function on a lattice. \emph{Operations Research} 26(2), 305--321.
\bibitem{milgrom} Milgrom, P.~R. (1981). Good news and bad news: representation theorems and applications. \emph{Bell Journal of Economics} 12(2), 380--391.
\bibitem{blackwell} Blackwell, D. (1953). Equivalent comparisons of experiments. \emph{Annals of Mathematical Statistics} 24(2), 265--272.
\end{thebibliography}

\end{document}
